\documentclass[10pt,a4paper]{amsart}
\usepackage{amsmath,amssymb,amsthm,graphicx}
\usepackage{verbatim}
\newenvironment{reference}{\comment}{\endcomment}
%\newenvironment{reference}{}{}
\newenvironment{slogan}{\comment}{\endcomment}
\newenvironment{history}{\comment}{\endcomment}

% For commutative diagrams we use Xy-pic
\usepackage[all]{xy}

% We use 2cell for 2-commutative diagrams.
\xyoption{2cell}
% Theorem environments.
%
\theoremstyle{plain}
\newtheorem{theorem}[subsection]{Theorem}
\newtheorem{proposition}[subsection]{Proposition}
\newtheorem{lemma}[subsection]{Lemma}

\theoremstyle{definition}
\newtheorem{definition}[subsection]{Definition}
\newtheorem{example}[subsection]{Example}
\newtheorem{exercise}[subsection]{Exercise}
\newtheorem{situation}[subsection]{Situation}

\theoremstyle{remark}
\newtheorem{remark}[subsection]{Remark}
\newtheorem{remarks}[subsection]{Remarks}

\numberwithin{equation}{subsection}

% Macros
%
\def\lim{\mathop{\mathrm{lim}}\nolimits}
\def\colim{\mathop{\mathrm{colim}}\nolimits}
\def\Spec{\mathop{\mathrm{Spec}}}
\def\Hom{\mathop{\mathrm{Hom}}\nolimits}
\def\Ext{\mathop{\mathrm{Ext}}\nolimits}
\def\SheafHom{\mathop{\mathcal{H}\!\mathit{om}}\nolimits}
\def\SheafExt{\mathop{\mathcal{E}\!\mathit{xt}}\nolimits}
\def\Sch{\mathit{Sch}}
\def\Mor{\mathop{\mathrm{Mor}}\nolimits}
\def\Ob{\mathop{\mathrm{Ob}}\nolimits}
\def\Sh{\mathop{\mathit{Sh}}\nolimits}
\def\NL{\mathop{N\!L}\nolimits}
\def\CH{\mathop{\mathrm{CH}}\nolimits}
\def\proetale{{pro\text{-}\acute{e}tale}}
\def\etale{{\acute{e}tale}}
\def\QCoh{\mathit{QCoh}}
\def\Ker{\mathop{\mathrm{Ker}}}
\def\Im{\mathop{\mathrm{Im}}}
\def\Coker{\mathop{\mathrm{Coker}}}
\def\Coim{\mathop{\mathrm{Coim}}}

% Boxtimes
%
\DeclareMathSymbol{\boxtimes}{\mathbin}{AMSa}{"02}

%
% Macros for moduli stacks/spaces
%
\def\QCohstack{\mathcal{QC}\!\mathit{oh}}
\def\Cohstack{\mathcal{C}\!\mathit{oh}}
\def\Spacesstack{\mathcal{S}\!\mathit{paces}}
\def\Quotfunctor{\mathrm{Quot}}
\def\Hilbfunctor{\mathrm{Hilb}}
\def\Curvesstack{\mathcal{C}\!\mathit{urves}}
\def\Polarizedstack{\mathcal{P}\!\mathit{olarized}}
\def\Complexesstack{\mathcal{C}\!\mathit{omplexes}}
% \Pic is the operator that assigns to X its picard group, usage \Pic(X)
% \Picardstack_{X/B} denotes the Picard stack of X over B
% \Picardfunctor_{X/B} denotes the Picard functor of X over B
\def\Pic{\mathop{\mathrm{Pic}}\nolimits}
\def\Picardstack{\mathcal{P}\!\mathit{ic}}
\def\Picardfunctor{\mathrm{Pic}}
\def\Deformationcategory{\mathcal{D}\!\mathit{ef}}
\usepackage{cleveref}
\crefname{originalreference}{original reference}{original references}
\Crefname{originalreference}{Original reference}{Original references}
\makeatletter
\newlabel{lemma-push-unit}{{0FHG}{source}{}{sourcepointer-1}{}}
\newlabel{lemma-push-unit@cref}{{[originalreference][0][]0FHG}{[1][0][]source}{}{}{}}
\makeatother
\crefname{originalreference}{original reference}{original references}
\Crefname{originalreference}{Original reference}{Original references}
\makeatletter
\newlabel{lemma-pushforward}{{03M9}{source}{}{sourcepointer-2}{}}
\newlabel{lemma-pushforward@cref}{{[originalreference][0][]03M9}{[1][0][]source}{}{}{}}
\makeatother
\crefname{originalreference}{original reference}{original references}
\Crefname{originalreference}{Original reference}{Original references}
\makeatletter
\newlabel{remark-replace-cup-product}{{0FHD}{source}{}{sourcepointer-3}{}}
\newlabel{remark-replace-cup-product@cref}{{[originalreference][0][]0FHD}{[1][0][]source}{}{}{}}
\makeatother
\title{Pairings of the Kunneth components of the class of the diagonal}
\author{The Stacks Project Authors}
\date{stacks-project:weil.tex:local-snapshot-20261012; GNU FDL 1.2 or later}
\begin{document}
\maketitle
\noindent\textit{Editorial note.} All delivered mathematical text is the author's own, copied from the pinned Stacks Project chapter file weil.tex; no statement or proof step was rewritten, reordered or removed.\par\medskip
\noindent\textit{Editorial note.} The counted proof is the author's complete argument: the identification of the class of the diagonal as the pushforward of the unit, the expansion through multiplicativity of the degree map, and the choice of dual class producing the displayed pairings.\par\medskip
\noindent\textit{Editorial note.} The label lemma-class-diagonal names the counted statement, and the proof environment immediately following it in the source is delivered verbatim as its proof.\par\medskip
\noindent\textit{Editorial note.} The counted statement is carried by the author's own lemma environment in the chapter source.\par\medskip
\noindent\textit{Editorial note.} Pre-open gate: python operations/agent-authoring/tools/precard\_gate.py --source weil.tex --label lemma-class-diagonal --start 2860 --end 2899 --local-id 28148 returned RESULT=OK (L28148 free, occupancy FREE, public interval clear). Pre-assembly tag self-check: canonical tag resolved and not present in reference/public-items.jsonl.\par\medskip
\noindent\textit{Editorial note.} The statement and proof are copied from the pinned chapter file; the cited inputs are kept as external pointers with their canonical Stacks tags.\par\medskip
\noindent\textit{Allowed standard prerequisites (editorial).}\begin{itemize}
\item The class of the diagonal as the pushforward of the unit. Same chapter, the lemma identifying the class of the diagonal with the pushforward of the unit, and the lemma expressing the degree map of a product through the pairings on the factors. Turns the pairing of the class of the diagonal with a product of classes into the pairing of those classes on the scheme.
\item Pullback along the diagonal is the cup product. Same chapter, the remark identifying the pullback along the diagonal with the cup product. Identifies the action of the diagonal with the multiplication in cohomology.
\item Behaviour of the degree map under products. Axiom (B)(b) on the multiplicativity of the degree map for products. Expands the pairing of a decomposable class with the Kunneth decomposition into a sum of products of pairings.
\end{itemize}\medskip
\ifdefined\hypertarget\hypertarget{sourcepointer-1}{}\fi
\noindent\textit{Source reference.} Original reference 0FHG: sources/stacks-weil/pinned-original-source.tex, source lines 2848--2848. Prerequisite pointer invoked at proof line(s) 2881; it names a standard fact of the same chapter that the counted argument uses, and every step of the counted proof built on it is retained verbatim in the delivered body.\par\medskip
\ifdefined\hypertarget\hypertarget{sourcepointer-2}{}\fi
\noindent\textit{Source reference.} Original reference 03M9: sources/stacks-weil/pinned-original-source.tex, source lines 2660--2660. Prerequisite pointer invoked at proof line(s) 2888; it names a standard fact of the same chapter that the counted argument uses, and every step of the counted proof built on it is retained verbatim in the delivered body.\par\medskip
\ifdefined\hypertarget\hypertarget{sourcepointer-3}{}\fi
\noindent\textit{Source reference.} Original reference 0FHD: sources/stacks-weil/pinned-original-source.tex, source lines 2742--2742. Prerequisite pointer invoked at proof line(s) 2880; it names a standard fact of the same chapter that the counted argument uses, and every step of the counted proof built on it is retained verbatim in the delivered body.\par\medskip
\begin{lemma}
\label{lemma-class-diagonal}
Assume given (D0), (D1), (D2), and (D3) satisfying (A), (B), and (C).
Let $X$ be a nonempty smooth projective scheme over $k$ equidimensional
of dimension $d$. Choose a basis $e_{i, j}, j = 1, \ldots, \beta_i$ of
$H^i(X)$ over $F$. Using K\"unneth write
$$
\gamma([\Delta]) =
\sum\nolimits_i
\sum\nolimits_j e_{i, j} \otimes e'_{2d - i , j}
\quad\text{in}\quad
\bigoplus\nolimits_i H^i(X) \otimes_F H^{2d - i}(X)(d)
$$
with $e'_{2d - i, j} \in H^{2d - i}(X)(d)$.
Then $\int_X e_{i, j} \cup e'_{2d - i, j'} = (-1)^i\delta_{jj'}$.
\end{lemma}

\begin{proof}
Recall that $\Delta^* : H^*(X \times X) \to H^*(X)$ is equal to the
cup product map $H^*(X) \otimes_F H^*(X) \to H^*(X)$, see
Remark \ref{remark-replace-cup-product}. On the other
hand, recall that $\gamma([\Delta]) = \Delta_*1$ (Lemma \ref{lemma-push-unit})
and hence
$$
\int_{X \times X} \gamma([\Delta]) \cup a \otimes b =
\int_{X \times X} \Delta_*1 \cup a \otimes b =
\int_X a \cup b
$$
by Lemma \ref{lemma-pushforward}.
On the other hand, we have
$$
\int_{X \times X} (\sum e_{i, j} \otimes e'_{2d -i , j}) \cup a \otimes b =
\sum (\int_X a \cup e_{i, j})(\int_X e'_{2d - i, j} \cup b)
$$
by axiom (B)(b); note that we made two switches of order so that the sign
for each term is $1$. Thus if we choose $a$ such that
$\int_X a \cup e_{i, j} = 1$ and all other pairings equal to zero, then
we conclude that $\int_X e'_{2d - i, j} \cup b = \int_X a \cup b$
for all $b$, i.e., $e'_{2d - i, j} = a$. This proves the lemma.
\end{proof}

\end{document}
